\documentclass[12pt]{article}
\usepackage[utf8]{inputenc}
\usepackage[T1]{fontenc}
\usepackage{hyperref}
\title{Yidiny Language Resources for K-12 Curriculum Development}
\author{Compiled by Claude Sonnet 5.5 (claude-sonnet-5-5) for Andr\'as Kornai}
\date{}
\begin{document}
\maketitle
\section{Introduction}
Yidiny (ISO 639-3 yii) is an endangered Pama-Nyungan language of the Yidinji people of north-east Queensland, in the Cairns and Tablelands regions. Wikipedia reports 52 speakers in the 2021 census. It is not known to be taught in schools. The search found academic descriptions but no school materials.
\section{Primary Educational Resources}
No primers, textbooks or school programmes for Yidiny were found in the search.
\section{Linguistic and Reference Materials}
R.\ M.\ W.\ Dixon, \emph{A Grammar of Yidi\~n} (Cambridge University Press, 1977), is the standard reference. Its contents include a lexicon section. Preface: \url{https://resolve.cambridge.org/core/books/grammar-of-yidin/preface/A7ED2595CD555C60CBC4122652DCD8DD}. Catalogue record: \url{https://ci.nii.ac.jp/ncid/BA01483322}

Dixon, \emph{Edible Gender, Mother-in-Law Style, and Other Grammatical Wonders} (2015), includes studies of Yidi\~n: \url{https://researchonline.jcu.edu.au/39196/}

No standalone dictionary was found in the search.
\section{Religious and Community Resources}
Background on the Yidiny people: \url{https://en.wikipedia.org/wiki/Yidiny_people}. No community language organisation was identified in the search.
\section{Digital and Audio Resources}
No audio or digital resources were found in the search.
\section{Curriculum Development Context}
Very few speakers remain and no school programme is known. A curriculum would need to be developed with the Yidinji community and would rely on Dixon's documentation. Archive holdings (for example AIATSIS) were not confirmed here and should be checked.
\end{document}